\subsection{构型依赖质量下的配对学习}
\label{sec:variable_mass_learning}

上述学习以常质量为基准，而式（36）的变质量项直接涉及不可分离动能。为在相同物理对象上检验这一分支，取可逆正则点变换
\[
 X_j=f(q_j)=q_j+q_j^3,\qquad
 Y_j=\frac{p_j}{f'(q_j)},\qquad j=1,2.
\]
该变换满足$\sum_jY_j\,dX_j=\sum_jp_j\,dq_j$，从而保持闭且非退化的正则二形式$\Omega=\sum_jdX_j\wedge dY_j=\sum_jdq_j\wedge dp_j$及其Jacobi结构；同一梁能量在新坐标中成为
\[
 \begin{aligned}
 H_{\rm vm}(q,p)&=\frac12p^TA(q)p+V(f(q)),\\
 A_{jj}(q)&=\frac{\nu_j}{(1+3q_j^2)^2}.
 \end{aligned}
\]
因此$M(q)=A(q)^{-1}$随构型变化，式（36）中的$\Delta_j=-\tfrac12\bar p_j^2\partial_{q_j}A_{jj}(\bar q)$一般非零。这是同一连续模型的坐标检验，不引入新的实体质量本构；非线性坐标下的中点离散也不等于把旧离散轨迹直接变换，故教师轨迹由新坐标中的式（36）重新求解。

四组网络沿用相同参数数目、更新预算、噪声和五个随机种子。G组固定的物理主部相应改为$H_{0,\rm vm}=\tfrac12\sum_j\nu_j[p_j^2/f'(q_j)^2+f(q_j)^2]$，避免把构型依赖动能误作未知残差。监督仍以一步差商配对新坐标中点；独立连续参考也在该坐标中积分。预测与参考统一映回原$(X,Y)$后计算状态误差和杆端RMSE，五项配对定义保持不变。网络预测采用所学Hamiltonian的通用中点方程；任意网络能量仍可能超出式（36）原有的质量--势能系数形式。

\begin{table*}[t]
\centering\small
\caption{变质量正则坐标检验的五项种子内配对差（预测步长$h_0/4$，杆端RMSE差，单位$\mu\mathrm m$）。}
\label{tab:variable_mass_paired}
\begin{tabular}{lrr}
\toprule
配对差 & 对同步长教师 & 对连续参考 \\
\midrule
$R-GR$ & $847.52\pm191.45$ & $843.41\pm191.90$ \\
$C-GC$ & $861.01\pm281.76$ & $856.68\pm281.50$ \\
$R-C$ & $-13.308\pm134.04$ & $-13.269\pm133.71$ \\
$GR-GC$ & $0.18377\pm0.06680$ & $-0.0001783\pm0.30789$ \\
$(GR-GC)-(R-C)$ & $13.492\pm134.01$ & $13.268\pm133.70$ \\
\bottomrule
\end{tabular}
\par\smallskip\begin{minipage}{0.96\textwidth}\footnotesize 每个种子内先平均8条独立测试轨迹的时间域RMSE，再构造五项配对差；表列5个种子的均值$\pm$样本标准差。正值表示前组误差较大。误差在原物理模态坐标中评价；该实验为同一梁的非线性正则坐标检验。\end{minipage}
\end{table*}


20次端到端训练及两档步长的40组预测均完成，隐式迭代均达到设定容差。式（36）矩阵更新与中点方程的独立代回差小于$1.8\times10^{-15}$，中点标签与Hamiltonian场的差小于$3.9\times10^{-14}$，该坐标光滑算例的最细观察阶约为$1.998$。在$h_0/4$下，GR和GC对同步长教师的杆端RMSE分别为$0.611\pm0.232$与$0.428\pm0.231\,\mu\mathrm m$，GC在五个种子均较低；对连续参考却分别为$4.901\pm0.782$和$4.901\pm0.525\,\mu\mathrm m$，GR--GC配对差相对种子波动接近零（表\ref{tab:variable_mass_paired}）。无网络物理主部的杆端误差为$4.338\pm0.439\,\mu\mathrm m$，完整状态误差约$0.006816$，也低于GR的$0.008028$与GC的$0.008056$。因此，非零$\Delta$下的数据生成和五项对照已经实现，而所测步长下网络尚未超越物理基线，复值拟合收益也未稳定转化为连续预测优势。
